\begin{table}[!htbp]
\centering
\caption{Robustness to regional spillovers, 9th Math: ATT by relative period}
\label{tab:ap_event_spillover_9th_mat}
\vspace{-0.15cm}
\scriptsize
\setlength{\tabcolsep}{4pt}
\renewcommand{\arraystretch}{0.92}
\begin{tabular}{lccc}
\toprule
Relative period & Main & Exclude exposed DE & Exclude exposed district-DE \\
\midrule
-4 & 0.060 (0.065) & 0.028 (0.069) & 0.023 (0.069) \\
-3 & 0.037 (0.040) & -0.094 (0.050) & -0.090 (0.048) \\
-2 & -0.003 (0.030) & -0.050 (0.039) & -0.062 (0.040) \\
-1 & 0.000 & 0.005 (0.033) & 0.000 (0.031) \\
0 & 0.310 (0.038) & 0.291 (0.045) & 0.327 (0.042) \\
1 & 0.442 (0.036) & 0.441 (0.051) & 0.440 (0.038) \\
2 & 0.602 (0.039) & 0.548 (0.071) & 0.610 (0.042) \\
3 & 0.751 (0.049) & 0.738 (0.092) & 0.764 (0.051) \\
4 & 0.901 (0.062) & 0.868 (0.078) & 0.934 (0.065) \\
\bottomrule
\end{tabular}
\begin{flushleft}
\footnotesize Note: each cell reports the ATT and standard error in parentheses, in standard deviations of annual proficiency. Period \(k=-1\) is normalized as the baseline and therefore appears without a standard error.
\end{flushleft}
\end{table}



